\documentclass[12pt,a4paper,oneside]{article}
\usepackage[brazil]{babel}
\usepackage[utf8]{inputenc}
\usepackage{olymp}


\newlength{\exmpwidouf}
\exmpwidinf=10cm
\exmpwidouf=6cm


\setcounter{problem}{2}

\begin{document}

\begin{problem}{Siglas}{siglas}{3}

Um acrônimo, ou mais comumente conhecido como sigla, é uma abreviação formada pelos componentes iniciais de uma palavra ou de uma frase. Estes componentes podem ser letras individuais, como em UFV (\underline{\textbf{U}}niversidade \underline{\textbf{F}}ederal de \underline{\textbf{V}}içosa), ou partes das palavras, como em ANVISA (\underline{\textbf{A}}gência \underline{\textbf{N}}acional de \underline{\textbf{Vi}}gilância \underline{\textbf{Sa}}nitária).

Na verdade, não existe um padrão muito rígido para a formação de siglas, e existem siglas que contém partes das palavras que não estão necessariamente no início delas, como em DPI (\underline{\textbf{D}}e\underline{\textbf{p}}artamento de \underline{\textbf{I}}nformática) ou em CCP (\underline{\textbf{C}}iência da \underline{\textbf{C}}om\underline{\textbf{p}}utação).

Faça um programa que, dada uma ou mais palavras com algumas letras destacadas em maiúsculas, gere a sigla correspondente que contém todas as letras destacadas. Veja os exemplos mais adiante.
%\Notes
%
%Observações, se tiver.

\InputFile

A entrada contém somente uma linha, com uma ou mais palavras. As palavras contêm somente caracteres simples maiúsculos e minúsculos, e são separadas por espaços. Restrições: a entrada contém no máximo 200 caracteres.


\OutputFile

Seu programa deve escrever a sigla correspondente às letras destacadas na entrada.

% \Notes
% Preste atenção aos tipos das variáveis, pois $F(90) \approx 2^{61}$.

\Example

\begin{example}
\exmp{
Universidade Federal de Vicosa
}{
UFV
}%
\end{example}

\begin{example}
\exmp{
DePartamento de Informatica
}{
DPI
}%
\end{example}

\begin{example}
\exmp{
Ciencia da ComPutacao
}{
CCP
}%
\end{example}


\begin{example}
\exmp{
Agencia Nacional de VIgilancia SAnitaria
}{
ANVISA
}%
\end{example}

\begin{example}
\exmp{
COronaVIrus Disease
}{
COVID
}%
\end{example}

%Comentários sobre o exemplo, se necessário.

\end{problem}

\end{document}
